% chapter2.tex
\section{相关理论基础}

\subsection{可变形物体物理建模理论}
可变形物体的物理建模是实现高真实感触觉交互的重要物理基础，其主要目的就是计算机虚拟环境中对软体、柔性材料、布料、人体器官等物体在接触、按压、滑动、揉捏、拖拽等动作过程中产生的形变、回弹、应力分布、能量耗散和接触情况做出准确、高效且稳定的数值计算。

刚体仿真只需要考虑位置、姿态、速度就可以，可变形物体要应对大变形、非线性力学、分布式接触、多质点耦合、时变材料属性等许多复杂的难题，属于触觉仿真领域最具有挑战性的核心技术之一。

\subsubsection{连续介质力学基础理论}
连续介质力学是所有可变形物体分析的宏观理论基础，它把物体看作无限密集、连续分布的质点群落，忽略原子、分子之间的空隙，只从场的角度来研究位移场、应变场、应力场和速度场之间的关系\cite{dai2025}。

根据仿真目标和材料特性，分别用两种经典的力学模型来建立本系统的模型。弹性力学模型适用于果冻、橡胶、弹性塑料等完全可恢复变形的材料，应力和应变是稳定的，和时间无关，没有滞后和蠕变，卸载后变形完全消失，没有残余形变。

粘弹性力学模型适合于肝脏、肌肉、皮肤、软组织等生物活体材料，这些材料既有弹性固体又有粘性流体的性质。存在应力松弛、蠕变现象，变形速度越大，感知刚度越高，存在明显的能量耗散、迟滞特性。

连续介质力学给后面离散化建模、本构方程设计、接触力计算提供了一个统一的物理依据。

\subsubsection{离散化建模方法}
计算机不能直接求解连续介质偏微分方程，因此需要将它离散化为可以计算的离散单元结构，该系统根据即时性、保真度、稳定性等要求，采用质点-弹簧模型与MuJoCo Flex改进后的柔性体模型相结合的混合建模方式。

\begin{enumerate}
    \item 质点--弹簧模型
    
    质点-弹簧模型（Mass-Spring Model，MSM）是即时交互领域里最成熟、使用最广泛的可变形体建模方法\cite{zhai2024}。该模型的基本结构是由质点（质量分布）、弹簧（弹性约束）和阻尼器（振动衰减）这三部分组成。
    
    \begin{figure}[H]
        \centering
        \includegraphics[width=0.8\textwidth]{figures/mass_spring.png}
        \caption{质点--弹簧模型拓扑结构原理图}
        \label{fig:mass_spring}
    \end{figure}
    
    其主要优势就是计算量很小，可以支持物理步长为1000Hz以上的计算；弹簧--阻尼结构直观，刚度、阻尼、质量参数容易调节；对大变形有较强的适应性，不容易产生网格翻转；与触觉刷新率（100～300Hz）天然匹配。系统通过调节刚度系数stiffness、阻尼系数damping和网格密度grid count来控制不同的材质。
    
    \item MuJoCo Flex柔性体实现机制
    
    MuJoCo具有高度改进的flexcomp柔性组件，这是本系统可变形物体的主要实现手段，底层对质点-弹簧模型进行数值优化和约束增强，
    \begin{enumerate}
        \item grid：规则三维网格，适合立方体、块状软体；
        \item ellipsoid：椭球结构，适合于球形、类球形的器官；
        \item gmsh：外部高精度网格，肝、肾等复杂模型用的较多。
    \end{enumerate}
    Flex组件用edge来定义质点的弹性结合，有接触约束、边界固定、自碰撞检测等功能，能很好地应对手指按压、滑动、揉捏等各种复杂的交互情况，不会出现穿透、飞散、震荡发散等问题。
\end{enumerate}

\subsubsection{材料本构关系}
本构关系就是反映应力和应变之间本质的数学联系，决定物理是否真实。

\begin{enumerate}
    \item 线性弹性本构适用于标准弹性材料：
    \begin{equation}
        \sigma = E \epsilon
    \end{equation}
    其中：$\sigma$为应力；$E$为弹性模量；$\epsilon$为应变。
    
    \item 粘弹性本构适用于生物软组织，增加粘性耗散项：
    \begin{equation}
        \sigma = E \epsilon + \eta \dot{\epsilon}
    \end{equation}
    其中：$\eta$为粘性系数；$\dot{\epsilon}$为应变速率。该模型能够准确还原“快按更硬、慢按更软”的真实生物组织触感。
\end{enumerate}

\subsection{触觉反馈与力触觉渲染理论}
触觉是人感知外部世界最直接、最精确的途径，在虚拟操作中，如果缺少触觉反馈，就会导致用力失控、定位不准、任务失败概率增大、沉浸感明显下降等问题。

触觉渲染的主要目的就是将物理引擎给出的接触信息转化为人体指尖可以感知到的压力、纹理、粗糙度、滑动、软硬等机械刺激信号，并且需要达到低时延、高刷新率、高保真度的效果。

\subsubsection{接触力计算模型}
\begin{enumerate}
    \item 基于Hertz的非线性接触模型
    
    指尖按压可变形物体时，接触力符合Hertz接触理论，更接近真实弹性体接触规律
    \begin{equation}
        F = k \delta^{\frac{3}{2}}
    \end{equation}
    其中，$F$是法向接触力，$k$是等效刚度，$\delta$是变形或者穿透深度。该模型有非线性刚度特性，即“越压越硬”，这是真实的触觉的主要特点。
    
    \item 自适应力反馈计算机制
    
    为了进一步提高真实感，系统建立多因素融合自适应力模型：
    \begin{enumerate}
        \item 基础力项：来自MuJoCo接触传感器；
        \item 深度补偿项：强化大变形下的刚度感知；
        \item 速度阻尼项：抑制高速冲击与震荡；
        \item 力限幅项：保证在人体安全感知范围内；
        \item 动态滤波项：消除毛刺与高频噪声。
    \end{enumerate}
\end{enumerate}

\subsubsection{力信号处理与平滑技术}
\begin{enumerate}
    \item 滑动平均滤波
    
    物理引擎输出的原始力信号含有高频噪声，如果直接输出就会引起刺痛、震颤等不良触觉\cite{ren2025}，因此系统采用滑动平均滤波的方法。
    \begin{equation}
        F_{\text{smooth}}(t) = \frac{1}{N} \sum_{i=0}^{N-1} F(t-i)
    \end{equation}
    窗口$N$为5的时候，平滑性以及响应速度达到了最佳的平衡。
    
    \item 冲击力检测与增强
    
    通过力的变化率识别快速撞击、拍打等事件：
    \begin{equation}
        r = \frac{dF}{dt}
    \end{equation}
    当变化率超过阈值的时候，就会认定是冲击，系统会自动加大反馈强度来改善触觉的层次感。
\end{enumerate}

\subsubsection{纹理触觉渲染机制}
\begin{enumerate}
    \item 动态纹理强度计算
    
    纹理是由滑动摩擦造成的高频振动产生的，压力、速度变化的时候强度也会跟着变化。
    \begin{equation}
        I_{\text{texture}} = I_0 \cdot (1 + \alpha F) \cdot (1 + \beta v)
    \end{equation}
    该式子说明压力增大时，花纹就更明显；滑动速度加快，振动强度就增大。各种材质所具有的触感是不一样的。
    
    \item 物体--纹理精准映射
    
    系统建立严格的材质纹理对应关系，达到高度逼真的材质感知效果：
    \begin{itemize}
        \item 果冻 $\rightarrow$ ProfiledRubberSlow（柔软弹性）
        \item 布料 $\rightarrow$ Smooth（光滑低摩擦）
        \item 海绵 $\rightarrow$ CrushedRock（粗糙多孔）
        \item 肝脏 $\rightarrow$ VenetianGranite（软组织细腻颗粒感）
    \end{itemize}
\end{enumerate}

\subsection{手部运动映射与逆运动学理论}
虚拟手是用户进入虚拟世界唯一的一个交互工具，如果虚拟手的运动是准确的、流畅的、自然的，那么操作起来就会感到舒适，任务完成得更加稳定，也会产生更强的沉浸感。

\begin{enumerate}
    \item 逆运动学基本原理
    
    逆运动学（IK）的核心问题就是：已知指尖目标位置 $\rightarrow$ 求解手指各个关节的角度。手指为多关节串联结构，传统的伪逆法（Pseudoinverse）有明显的不足：奇异点附近关节剧烈抖动，关节容易超出生理范围，速度突变造成动作僵硬。
    
    \item 阻尼最小二乘法（DLS）
    
    本系统使用阻尼最小二乘法DLS作为主要的IK算法，加入阻尼项来抑制奇异值发散
    \begin{equation}
        \dot{\theta} = J^T (J J^T + \lambda^2 I)^{-1} \dot{x}
    \end{equation}
    其中：$J$为雅可比矩阵；$\lambda$为阻尼系数；$\dot{x}$为末端速度；$\dot{\theta}$为关节速度。
    
    \begin{figure}[H]
        \centering
        \includegraphics[width=0.6\textwidth]{figures/ik_flow.png}
        \caption{DLS 逆运动学求解流程图}
        \label{fig:ik_flow}
    \end{figure}
    
    \item 动态阻尼调整策略
    
    系统要实现运动快、稳的目标，就要根据当前手指的姿态来自主调节阻尼的大小，当手指处于正常、开阔的运动范围内时，就用较小的阻尼来加快响应速度，使手指更加灵活。一旦手指接近奇异姿态，即容易产生抖动的时候，系统就会自动提高阻尼值，使运动变慢、更加稳定，在灵活应对和稳固不动之间找到最佳的平衡点。
    
    \item 关节限位与生理约束
    
    系统为了使动作更自然，设置了严格的约束条件，即各个关节的转动幅度不能超出人体手指实际的生理范围，不能出现反向弯曲、过度弯曲等不符合现实的动作，并且要防止手指互相交叉、碰撞，使虚拟手始终处于一个恰当、自然、实用的状态。
\end{enumerate}

\subsection{MuJoCo 物理引擎核心理论}
MuJoCo是目前全世界最符合触觉交互、机器人即时仿真物理引擎、高精度、高稳定、低时延、接触求解能力强的机器人仿真软件\cite{andras2023}。

\begin{enumerate}
    \item 多刚体与柔性体动力学
    
    MuJoCo依靠成熟的动力学理论，使用高效的递归计算方法，可以立即算出机器人手、虚拟手指、可变形物体的运动和受力情况，无论是刚体的抓取还是软体的挤压变形，都能给出稳定的、准确的、接近现实的物理效果。
    
    \begin{figure}[H]
        \centering
        \includegraphics[width=0.6\textwidth]{figures/mujoco_dynamics.png}
        \caption{MuJoCo 动力学求解循环图}
        \label{fig:mujoco_dynamics}
    \end{figure}
    
    \item 时间积分方法
    
    MuJoCo如果长时间运行而不会崩溃、发散，就需要使用稳定性较好的隐式积分法，该方法更适合于大变形、高刚度的物体，即使反复按压、揉捏软体，也不会出现模型乱飞、爆炸、错位等情况。
    
    \item 约束求解机制
    
    MuJoCo使用专门的约束求解算法，统一处理运动限制、碰撞阻挡、接触摩擦等各方面的问题，可以稳定地识别出手指和软体的接触情况，准确地计算出接触力，不会有物体相互穿透现象发生，这也是整个仿真系统稳定运行的重要保证。
    
    \begin{figure}[H]
        \centering
        \includegraphics[width=0.7\textwidth]{figures/constraint_solver.png}
        \caption{MuJoCo 约束求解流程示意图}
        \label{fig:constraint_solver}
    \end{figure}
\end{enumerate}

\subsection{实时仿真与多线程同步理论}
触觉交互属于一种对即时性、同步性要求很高的系统，一旦出现卡顿、延误或者错位的现象，就会直接影响到沉浸感的营造。

\begin{enumerate}
    \item 实时性核心指标
    
    系统需要满足严格的即时指标，物理仿真处于高频状态时，触觉反馈的更新频率要高，画面渲染要流畅，而且从用户动作到触觉输出的总体延误要控制在很低的水平，否则用户就会感觉到“手先动，触感后到”的现象，体验会大大降低。
    
    \item 三线程分离架构
    
    系统采用仿真、触觉、渲染完全分离并行架构：
    \begin{itemize}
        \item 仿真线程：专门负责物理计算和接触检测；
        \item 触觉线程：专门处理力反馈、滤波和设备驱动；
        \item 渲染线程：专门负责画面显示和VR视角更新。
    \end{itemize}
    三个线程同时运行，不会互相影响，可以将整个系统的延迟降到最低，使系统流畅度达到较高水平。
    
    \item 系统延迟构成与优化
    
    整个系统存在一定的延迟，它的源头是传感器采集数据、算法执行计算、画面完成渲染、触觉设备响应操作等环节。为了尽可能减小延迟，该系统采取了运动预测、时间戳补偿、提高线程优先级等许多改进措施，把总体延迟保持在人类几乎察觉不到的程度上，达到即时互动的效果。
\end{enumerate}

\subsection{触觉设备接口与驱动理论}
本系统可以使用WEART TouchDIVER触觉手套，该手套是实现多模态触觉反馈的重要的硬件。

\begin{enumerate}
    \item WEART工作原理
    
    该触觉手套通过指尖的振动、微小压力刺激等方式向用户传递压力、纹理、粗糙、滑动等各方面的信息，可以区分出软、硬、滑、糙等各种材质的特点，是实现高真实感虚拟触摸的重要硬件支持。
    
    \item 自动校准流程
    
    系统启动后要完成完整的自动校准过程才能保证反馈准确，该过程分为四步：第一步确定手指的初始零位，第二步创建力的输出映射关系，第三步校准纹理强度，第四步展开验证检测，从而保证每次触觉输出都是准确、稳定的、可重复的。
\end{enumerate}